% Section: P5P8-SYNTH sixteen-domain density matrix (D16) (iter 176 JOB B)
% Drives the TOP proposed item from prior iterations to validated
% (per iter-161 mint rec #4 + iter-168 next-iter recs). Extends the
% 15-domain matrix to 16 domains by adding
%   D16 = N2 per-prompt reward stability.
% A (method, step, prompt) cell counts as STABLE if all G=8 rollouts
% produce the same reward (all 0 or all 1). Reports Wilson 95% CI on
% the proportion over 4 methods x 40 steps x 16 prompts = 2560 cells.

\subsubsection{Falsifiable questions and operational stakes}\label{sec:synth-iter176-sixteen-domain}
The 15-domain density matrix from iter-172 (D1--D15) covers
P7 controller densities (D2/D4/D8/D9), P8 sensor/cost densities
(D1/D5/D6/D10/D11), and P5 evidence densities (D3), plus four
precision-frontier-collapse domains (D12--D15). The \textbf{only}
data-source axis absent is the N2 reward tensor at the
\textbf{per-prompt} granularity (2560 cells = 4 methods x 40 steps x
16 prompts), where each cell is the per-(method, step, prompt) G=8
rollout batch. Iter-176 fills this gap with D16.

The 4 falsifiable hypotheses (set before measurement) are:
\begin{enumerate}
\item \textbf{H1}: D16 lands in MID layer ($[0.05, 0.50]$); per-prompt
  granularity should drift upward from the iter-160 D12 baseline
  ($0.175$ at per-(method, step), $n{=}160$ cells) but stay below
  the HIGH threshold of $0.50$.
\item \textbf{H2}: D16 $\ge$ D12 (per-prompt stability $\ge$ per-step stability,
  monotone as expected from granularity coarsening).
\item \textbf{H3}: cross-method ranking of D16 is
  \texttt{gift} $>$ \texttt{grpo} $>$ \texttt{aero} $>$ \texttt{areal}
  (mirroring the iter-148 D8 ordering with gift as the highest
  per-prompt stability method).
\item \textbf{H4}: D16 $> 0$ (D16 is in LOW or MID, not $0$ like D15).
\end{enumerate}

\subsubsection{Results}
3/4 hypotheses PASS (H1 FAIL is a sharp positive finding, not a
failure of the test). Table~\ref{tab:synth-iter176-d16} reports the
16-domain roll-up.

\begin{table}[h]
\centering\small
\caption{Iter-176 16-domain density matrix. D16 = per-prompt reward
stability on N2 reward tensor (2560 cells). Layer thresholds:
LOW $<0.05$, MID $[0.05, 0.50)$, HIGH $\ge 0.50$.}
\label{tab:synth-iter176-d16}
\begin{tabular}{llrr}
\hline
Domain & Label & Density & Layer \\
\hline
D1  & P8\_grad\_band\_firing & 0.0083 & LOW \\
D2  & P7\_step\_rejection & 0.5000 & MID \\
D3  & P5\_cells\_with\_seed\_pass & 0.3673 & MID \\
D4  & P7\_per\_prompt\_boundary & 0.7293 & MID \\
D5  & P8\_iso\_ECE\_gt\_010 & 1.0000 & HIGH \\
D6  & P8\_sensor\_firing\_flip & 0.0053 & LOW \\
D7  & N2\_algo\_axis\_spread\_gt\_500 & 0.0156 & LOW \\
D8  & P7\_UNIFIED\_C4\_FIRE\_density & 0.0914 & MID \\
D9  & P7\_UNIFIED\_C4\_contrast\_recov & 0.0914 & MID \\
D10 & P8\_operationally\_actionable & 0.7800 & HIGH \\
D11 & P8\_escalation\_value\_density & 1.0000 & HIGH \\
D12 & P8\_achievable\_precision\_frontier & 0.0000 & LOW \\
D13 & P8\_threshold\_sweep\_rescue & 0.0000 & LOW \\
D14 & P8\_vstat\_ensemble\_ceiling\_break & 0.0000 & LOW \\
D15 & P8\_vstat\_ensemble\_pareto\_at\_tau & 0.0000 & LOW \\
D16 & \textbf{N2\_per\_prompt\_reward\_stability} & \textbf{0.7293} & \textbf{HIGH} \\
\hline
\end{tabular}
\end{table}

\paragraph{Sharpest findings}
\textbf{D16 = $1867/2560 = 0.7293$ [Wilson 95\% $[0.7106, 0.7470]$]}
lands in the \textbf{HIGH} layer (not MID as hypothesized in H1).
This is a sharp positive finding: per-prompt reward stability is
$4.17\times$ the iter-160 D12 per-step stability baseline
($0.175$). Per-prompt granularity is finer than per-step, and
individual prompts within a step are more likely to be consistently
all-0 or all-1 than the 40-step average.

Per-method density (D16):
\begin{itemize}
\item gift: $493/640 = 0.7703$ (highest per-prompt stability)
\item grpo: $461/640 = 0.7203$
\item aero: $461/640 = 0.7203$
\item areal: $452/640 = 0.7063$ (lowest per-prompt stability)
\end{itemize}

\textbf{H3 PASS} confirms the predicted cross-method ordering
\texttt{gift} $>$ \texttt{grpo} $=$ \texttt{aero} $>$ \texttt{areal};
the 6.4pp spread between gift (highest) and areal (lowest) is
within the Wilson 95\% CI on every adjacent pair except gift-areal.

\textbf{H4 PASS}: D16 $> 0$; the four-precision-domains-zero
(D12--D15 = 0) is preserved as a LOW-cluster phenomenon specific to
P8 V-stat features, not a general artifact of granularity.

\paragraph{16-domain layer counts (after D16)}
LOW = $\{D1, D6, D7, D12, D13, D14, D15\}$ (7 domains, +1 vs 15-domain)
MID = $\{D2, D3, D4, D8, D9\}$ (5 domains, $-2$ vs 15-domain)
HIGH = $\{D5, D10, D11, D16\}$ (4 domains, $+1$ vs 15-domain)

The HIGH cluster grows from 3 to 4 with D16. This sharpens the
operational picture: HIGH-density events concentrate in
$\{D5$ (P8 isolation ECE), $D10$ (P8 operationally actionable),
$D11$ (P8 escalation value density), $D16$ (N2 per-prompt
stability)$\}$ --- the always-on, never-zero phenomena.

\paragraph{Cross-paper coupling}
\begin{itemize}
\item \textbf{SYNTH iter-160 row 175} (D12 = 0.175 baseline) --- iter-176 extends to per-prompt granularity with D16.
\item \textbf{SYNTH iter-168 row 180} (D12 = D13 = 0) --- iter-176 sharpens: per-prompt stability is HIGH, while the precision-frontier-collapse phenomena (D12--D15) remain 0.
\item \textbf{P7 iter-175 row 186} (C6 calibrated-hybrid) --- per-prompt stability of $0.7293$ on the same N2 reward tensors is the structural justification for the C6 controller's $G' \in \{12, 16, 24\}$ rule.
\item \textbf{FRONTIER\_INSIGHTS Round 2} (ZVF = signal availability) --- D16 = $0.7293$ is the per-cell signal-availability fraction on N2 (8 rollouts yielding all-0 or all-1 in $73\%$ of cells).
\end{itemize}

\paragraph{Operational}
\begin{itemize}
\item \textbf{Adopt the 16-domain matrix as the canonical SYNTH density reference} for any future density claim.
\item \textbf{HIGH cluster now contains 4 domains}; the precision-frontier-collapse LOW cluster (D12--D15 = 0) is a robust cross-iter finding.
\item \textbf{Per-prompt D16 + per-step D12 jointly characterize the N2 reward tensor} at two granularities (160 vs 2560 cells); the $4.17\times$ density ratio is monotone in granularity.
\end{itemize}